\section{Autoregressive image modeling：统一接口为何没有赢得所有任务}

如果 image 可以离散成 tokens，最自然的想法是沿用 GPT：text-to-image、image-to-text、image completion 都只是不同 token 排列与 mask。这条路线具有接口统一的美感，但 high-resolution sequence length、tokenizer information loss 和 sequential decoding 让它在生成质量与速度上面对 diffusion 的强竞争。

\subsection{CogView/DALL-E/Parti：先把图像变成离散序列}

设 image tokenizer 把图像 $I$ 编成 codes $z=(z_1,\ldots,z_m)$，文本 tokens 为 $x$。text-to-image 的 factorization 为
\[
p(z\mid x)=\prod_{i=1}^{m}p(z_i\mid x,z_{<i}).
\]
训练只需把 text 与 image tokens 拼成一个 sequence，但 image token 数可能达到数千，生成必须串行走完全部位置。

\lecturefigure{slide-22-autoregressive-text-to-image.jpg}{Autoregressive text-to-image：image tokenizer 加 GPT sequence}{Stanford Online 官方录像 00:46:47}

\begin{importantbox}{Image tokenizer 同时提供压缩和损失}
VQ tokenizer 把连续像素压成有限 codebook，显著缩短序列；但细文字、纹理、几何和颜色精度可能在编码阶段丢失。后面的 Transformer 无法恢复 tokenizer 从未保留的信息。
\end{importantbox}

\teachervoice{00:45:00--00:49:32，讲者把 CogView 与 DALL-E/Parti 放在同一简单 recipe：先离散化，再 autoregress。简单是优点，也是速度和细节瓶颈的来源。}

\subsection{Universal modeling：能做所有方向，不等于每个方向最好}

同一个 sequence model 可以通过排列和 mask 表达多种任务：
\[
\begin{aligned}
\text{text}\rightarrow\text{image}:&\ [x;z],\\
\text{image}\rightarrow\text{text}:&\ [z;x],\\
\text{masked image}:&\ [z_{\text{visible}};z_{\text{target}}].
\end{aligned}
\]
这确实统一了 API，但 image understanding 可能不如保留连续 features 的 VLM，image generation 又可能不如 diffusion。

\lecturefigure{slide-23-universal-modeling.jpg}{Universal vision-language modeling：排列统一与性能折中}{Stanford Online 官方录像 00:49:32}

\begin{warningbox}{统一 objective 与统一最优解是两件事}
一个模型能表达多任务 likelihood，不代表有限参数、有限数据和有限 compute 下会同时达到各 specialized systems 的 Pareto frontier。共享表示会带来 transfer，也会带来 interference 和不合适的 bottleneck。
\end{warningbox}

\teachervoice{00:49:32--00:52:11，讲者对自己参与的 universal modeling 路线给出很克制的评价：它实现了统一，但 image generation 不如 diffusion，image understanding 也不如连续视觉特征 VLM。}

\subsection{本章小结}

Autoregressive image modeling 最大优势是接口统一和成熟的 language-model machinery；最大问题是离散 tokenizer 损失、长序列串行生成和二维空间关系的线性化。它并非不可能生成好图，而是 practical frontier 转向了更适配并行图像计算的 diffusion。

